\chapter{Adenomyosis}

\section*{Definition and Overview}
Adenomyosis is a gynaecological condition in which tissue similar to the lining of the womb (endometrium) grows into the muscular wall of the womb (myometrium). The misplaced tissue thickens and bleeds with each menstrual cycle, making the womb enlarged, tender, and painful.

It most often affects women in their 30s and 40s and is a common cause of heavy, painful periods. It is closely related to endometriosis, in which similar tissue grows outside the womb, and the two conditions frequently occur together.

\section*{Epidemiology}
Adenomyosis is common and probably under-diagnosed. Studies using ultrasound and MRI suggest it may affect around 1 in 5 women, and it is found even more often in women undergoing hysterectomy.

It is most commonly diagnosed between the ages of 30 and 50 and is more likely in women who have had children. It often occurs together with endometriosis and uterine fibroids, which can make the clinical picture complex.

\section*{Aetiology and Pathophysiology}
The exact cause is not fully understood. The leading theory is that the boundary between the lining of the womb and its muscle wall is damaged, allowing lining tissue to push into the muscle. This damage may result from childbirth, caesarean section, or surgery on the womb.

The condition is driven by the hormone oestrogen. Each month the misplaced tissue responds to the cycle hormones, thickens, and bleeds, but the blood cannot drain normally, so the womb becomes swollen, heavy, and painful. Because the process depends on cyclical hormones, symptoms typically ease after menopause.

\section*{Clinical Features}
Symptoms vary considerably, and some women have none at all. Common features include:

\begin{itemize}
    \item \textbf{Heavy bleeding:} heavy, prolonged, or painful menstrual bleeding
    \item \textbf{Period pain:} severe cramping (dysmenorrhoea) that can last beyond the period
    \item \textbf{Pelvic pain:} chronic pelvic pain with a feeling of heaviness or pressure
    \item \textbf{Painful intercourse:} discomfort during or after sex
    \item \textbf{Enlarged womb:} bloating or a visibly enlarged womb
    \item \textbf{Fertility:} difficulty becoming pregnant in some women
    \item \textbf{Examination findings:} a tender, bulky womb on pelvic examination
\end{itemize}

\section*{Differential Diagnosis}
Other conditions cause similar symptoms and must be distinguished:

\begin{itemize}
    \item Uterine fibroids
    \item Endometriosis
    \item Dysfunctional (hormone-related) uterine bleeding
    \item Endometrial polyps
    \item Pelvic inflammatory disease
    \item Endometrial cancer (especially after menopause)
\end{itemize}

\section*{When to Seek Medical Care}
See a doctor if periods are becoming heavier or more painful, if they are affecting daily life, or if there is pelvic pain, pain during sex, or difficulty conceiving. Early assessment allows symptoms to be managed and other, more serious conditions to be ruled out.

\textbf{Call 000 (or local emergency services) for:}
\begin{itemize}
    \item Very heavy bleeding soaking pads or tampons rapidly
    \item Severe or sudden pelvic pain
    \item Fever with pelvic pain (possible infection)
    \item Dizziness, fainting, or signs of severe blood loss
    \item Any bleeding after menopause
\end{itemize}

\section*{Diagnosis and Investigations}
\begin{itemize}
    \item \textbf{History:} description of periods, bleeding patterns, pain, and fertility concerns
    \item \textbf{Pelvic examination:} the womb may feel enlarged and tender
    \item \textbf{Transvaginal ultrasound:} the usual first test, which can show thickening of the womb wall
    \item \textbf{MRI:} gives detailed images and is used when ultrasound is not conclusive
    \item \textbf{Blood tests:} full blood count and iron studies to check for anaemia
    \item \textbf{Endometrial biopsy or hysteroscopy:} occasionally used to rule out other conditions
\end{itemize}

\section*{Management}
Treatment depends on the symptoms, the desire for pregnancy, and the woman's age:

\begin{itemize}
    \item \textbf{Pain relief:} anti-inflammatory medicines such as ibuprofen
    \item \textbf{Tranexamic acid:} a medicine that reduces heavy bleeding
    \item \textbf{Hormonal treatment:} the combined oral contraceptive pill, the Mirena (levonorgestrel) IUD, progestogens, or medicines that lower oestrogen
    \item \textbf{Endometrial ablation:} a procedure to destroy the womb lining, which may help some women
    \item \textbf{Hysterectomy:} removal of the womb, the only treatment that cures adenomyosis
    \item \textbf{Fertility treatment:} specialist support for women trying to conceive
    \item \textbf{Supportive care:} counselling and pain-management strategies for chronic symptoms
\end{itemize}

\section*{Complications}
\begin{itemize}
    \item Iron-deficiency anaemia from heavy periods
    \item Chronic pelvic pain affecting work and daily life
    \item Reduced fertility and complications in pregnancy in some women
    \item Side effects of treatments, including hormonal effects
    \item Psychological distress, anxiety, and reduced quality of life
\end{itemize}

\section*{Prognosis and Outcomes}
Adenomyosis is a chronic condition that often improves after menopause, when hormone levels fall. Many women manage symptoms well with medicines and pain relief.

For women who choose hysterectomy, symptoms are usually cured, but the operation ends fertility and has its own recovery period and risks. With the right treatment and follow-up, most women find their symptoms substantially improved and their quality of life restored.

\section*{Prevention and Self-Care}
\begin{itemize}
    \item Seek treatment early for painful or heavy periods
    \item Use effective pain relief, taken early in the period
    \item Consider hormonal treatments such as the Mirena IUD if suitable
    \item Eat an iron-rich diet and treat anaemia if it develops
    \item Use heat, rest, and gentle exercise to manage pain
    \item Keep a symptom diary to help your doctor tailor treatment
    \item Discuss fertility plans early with your doctor if you are trying to conceive
\end{itemize}

\section*{Sources and Further Reading}
The following organisations provide reliable information on adenomyosis for patients and health professionals.

\begin{itemize}
    \item NHS. \textit{Information on adenomyosis, symptoms, and treatment.} https://www.nhs.uk/conditions/
    \item Mayo Clinic. \textit{Guide to adenomyosis, diagnosis, and management.} https://www.mayoclinic.org/
    \item MedlinePlus. \textit{Health information on adenomyosis and menstrual disorders.} https://medlineplus.gov/
    \item MSD Manual. \textit{Professional information on gynaecological disorders.} https://www.msdmanuals.com/
    \item NICE. \textit{Clinical guidance on heavy menstrual bleeding and gynaecological care.} https://www.nice.org.uk/
\end{itemize}
